\section{Tag 8 — Reed-Solomon-Decoder}

\subsection*{Bleistiftübung 1 — Syndrome berechnen}

Empfangenes Polynom (Auslösch-Position auf 0):
$r = (6, 1, 0, 5, 2)$, also $r(x) = 6 + x + 0 x^2 + 5 x^3 + 2 x^4$.

(a) $S_1 = r(\alpha) = r(2)$. Horner-Schema, höchster Koeffizient
zuerst, in GF($2^3$):
\begin{align*}
  &\text{start} = 2 \\
  &2 \cdot 2 \oplus 5 = 4 \oplus 5 = 1 \\
  &1 \cdot 2 \oplus 0 = 2 \\
  &2 \cdot 2 \oplus 1 = 4 \oplus 1 = 5 \\
  &5 \cdot 2 \oplus 6 = 1 \oplus 6 = 7
\end{align*}

(Hilfsrechnungen aus der GF($2^3$)-Multiplikationstafel:
$2 \cdot 2 = 4$, $5 \cdot 2 = 1$.)

\textbf{$S_1 = 7$.}

(b) $S_2 = r(\alpha^2) = r(4)$:
\begin{align*}
  &\text{start} = 2 \\
  &2 \cdot 4 \oplus 5 = 3 \oplus 5 = 6 \\
  &6 \cdot 4 \oplus 0 = 5 \\
  &5 \cdot 4 \oplus 1 = 2 \oplus 1 = 3 \\
  &3 \cdot 4 \oplus 6 = 7 \oplus 6 = 1
\end{align*}

(Hilfsrechnungen: $2 \cdot 4 = 3$, $6 \cdot 4 = 5$, $5 \cdot 4 = 2$,
$3 \cdot 4 = 7$.)

\textbf{$S_2 = 1$.}

(c) Beide Syndrome $\ne 0$ — es gab Fehler. (Was wir wussten, weil
wir selbst die Auslöschung erzeugt haben.)

\subsection*{Bleistiftübung 2 — Auslöschung korrigieren}

Position der Auslöschung: $j = 2$. Aus
$S_1 = e_j \cdot \alpha^j = e_2 \cdot \alpha^2$ folgt
\[
  e_2 = S_1 \cdot \alpha^{-2} = 7 \cdot \alpha^{-2}.
\]

(a)/(b) Aus Tag 7 wissen wir: $\alpha^7 = 1$, also $\alpha^{-2} =
\alpha^5 = 7$. Damit
\[
  e_2 = 7 \cdot 7.
\]
Aus der Multiplikationstafel: $7 \cdot 7 = 3$ (denn
$(x^2+x+1)^2 = x^4+x^2+1 = (x^2+x)+x^2+1 = x+1 = 3$).

\textbf{$e_2 = 3$.}

(c) Verifikation mit $S_2 = e_2 \cdot \alpha^4 = 3 \cdot \alpha^4 =
3 \cdot 6$. Aus der Tafel: $3 \cdot 6 = 1$ — und das ist genau das
$S_2$ aus Bleistiftübung 1. \checkmark

(d) Korrigiertes Symbol: $r_2 \oplus e_2 = 0 \oplus 3 = 3$. Das
ursprüngliche Codewort war $c = (6, 1, \mathbf{3}, 5, 2)$ — passt.

\subsection*{Aufgabe 1.1 (Python)}

Bei \emph{einem} gestörten Symbol bekommt man genau \emph{zwei}
Syndrome ungleich null (für $t = 4$ Prüfsymbole sogar vier — alle
sind nicht-null). Die Werte bilden eine geometrische Folge in
$\alpha^j$ mit Vorfaktor dem Fehlerwert.

Bei \emph{zwei} gestörten Symbolen bekommt man immer noch alle vier
Syndrome ungleich null, aber die Werte folgen keiner einzelnen
geometrischen Folge mehr — sie sind die Summe von zwei.

Daran kann man bei einer Inspektion ablesen, ob ein einzelner Fehler
oder mehrere vorliegen.

\subsection*{Aufgabe 2.1 (Python) — Bis an die Grenze}

Mit 4 Prüfsymbolen funktioniert die Korrektur für 1, 2, 3, 4
Auslöschungen. Bei 5 wirft \texttt{auslosung\_korrigieren} einen
\texttt{ValueError} („zu viele Auslöschungen“) — das System wäre
unterbestimmt (mehr Unbekannte als Gleichungen).

\subsection*{Aufgabe 2.2 (Python) — Falsche Position}

Wenn man eine Position als „Auslöschung“ angibt, an der eigentlich
kein Fehler ist, bekommt das System dort eine Unbekannte mit Wert 0.
Das System bleibt lösbar (sofern $v \le t$); die Lösung wird aber
für diese Position $e_j = 0$ liefern und den Wert nicht ändern. Mit
anderen Worten: einfach eine zusätzliche Auslöschung anzugeben „kostet
ein Prüfsymbol“ — man verbraucht eine Korrekturkapazität, ohne real
einen Fehler zu beheben.

\subsection*{Aufgaben 3.1 / 3.2 (Python, reedsolo)}

\textbf{3.1:} bei drei Symbol-Fehlern wirft \texttt{rsc.decode(\dots)}
eine Ausnahme vom Typ \texttt{reedsolo.ReedSolomonError} („Could not
locate error“ o.\,ä.). Genau das verlangt die Theorie: bei $n-k = 4$
Prüfsymbolen sind $\lfloor 4/2 \rfloor = 2$ Symbol-Fehler die
Garantie. Drei Fehler liegen jenseits.

\textbf{3.2:} mit dem Aufruf
\begin{verbatim}
korrigiert, _, _ = rsc.decode(daten, erase_pos=[1, 3, 5, 7])
\end{verbatim}
gibt man bis zu $n - k = 4$ Auslöschungspositionen mit. Der Decoder
nutzt diese als Hinweis und schafft dann auch 4 Auslöschungen oder
1 echter Fehler + 2 Auslöschungen usw. Generelle Regel:
$2 \cdot |\text{Fehler}| + |\text{Auslöschungen}| \le n - k$.

\subsection*{Antworten zu den drei Reflexionsfragen}

\begin{enumerate}
\item Wenn $v > t$, hat das Vandermonde-System mehr Unbekannte als
  Gleichungen — es gibt dann \emph{viele} mögliche Fehlerwert-Vektoren,
  die alle die Syndrome erklären. Der Decoder kann nicht eindeutig
  korrigieren. Anschaulich: wir haben zu viele „Löcher“ und zu wenig
  Information.

\item Ein zerkratztes Versandetikett liefert dem Scanner natürliche
  Auslöschungen: an den schwarzen Schlieren \emph{sieht der Scanner
  einfach kein Modul mehr}. Der Scanner weiß also: „hier fehlt
  Information“ — das ist eine Auslöschung mit bekannter Position.
  Genau dafür ist Reed-Solomon mit Auslöschungs-Korrektur gemacht.

\item Pragmatisch: laut „nicht lesbar“ und um einen erneuten Scan
  oder eine Kamera-Ausrichtung bitten. Stille Fehlversuche sind
  gefährlich: bei einer Postsendung könnte das Etikett
  fehlerhaft \emph{erfolgreich} dekodiert werden und das Paket zur
  falschen Adresse gehen. Im Zweifel ist Erkennung wichtiger als
  Korrektur — genau die Lehre aus CRC vs.\ Hamming, die wir an
  Tag 4 hatten.
\end{enumerate}
